\documentclass[aspectratio=169]{beamer}

\usepackage[utf8]{inputenc}
\usepackage[T1]{fontenc}
\usepackage{amsmath,amssymb}
\usepackage{bm}

\usetheme{Madrid}

\title{PhD Thesis Review and Synthesis}
\subtitle{Quantum Computing: From Physical Implementation to Practical Applications}
\author{Morten HJ}
\date{April 28}

\begin{document}

\frame{\titlepage}

%================================================
\section{Overall Assessment}
%================================================

\begin{frame}{Overall Assessment}
This thesis addresses a highly relevant and ambitious program in quantum computing:
\[
\text{Hardware} \rightarrow \text{Control} \rightarrow \text{Compilation} \rightarrow \text{Applications}
\]

\medskip

Strengths:
\begin{itemize}
    \item Strong physical motivation (electrons on helium)
    \item Relevant algorithmic contributions (variational compiling)
    \item NISQ-aware methodology
    \item Clear interdisciplinary direction
\end{itemize}

\medskip

Main issue:
\[
\text{Chapters are strong individually, but may need a  more unified thesis-level narrative/story/red thread}
\]
\end{frame}

%================================================
\section{Chapter-by-Chapter Review}
%================================================

\begin{frame}{Chapter 2: Theory}
\textbf{Strengths}
\begin{itemize}
    \item Clear introduction to quantum computing fundamentals
    \item Includes density matrices, entropy, and noise
\end{itemize}

\medskip

\textbf{Weaknesses}
\begin{itemize}
    \item Try to tailor more  to thesis problems
    \item Missing explicit links to later chapters
\end{itemize}

\medskip

\textbf{Required improvement}
\[
\rho_A = \mathrm{Tr}_B(\rho_{AB}), \quad
S(\rho) = -\mathrm{Tr}(\rho \log \rho)
\]

Explicit formalism should be connected to:
\begin{itemize}
    \item noise in compiling
    \item subsystem computation
\end{itemize}
\end{frame}

\begin{frame}{Chapter 3: Quantum Engineering}
\textbf{Strengths}
\begin{itemize}
    \item Frontier platform: electrons on helium
    \item Realistic Hamiltonian modeling
    \item Strong computational pipeline
\end{itemize}

\medskip

\textbf{Weaknesses}
\begin{itemize}
    \item Results described but not analyzed, link more to papers
    \item Could add  comparison to other qubit platforms
\end{itemize}

\medskip


Explicit structure and scaling discussion could be included.
\end{frame}

\begin{frame}{Chapter 4: Quantum Compiling}
\textbf{Strengths}
\begin{itemize}
    \item Strong theoretical contribution
    \item Recursive VQC idea is important
    \item Addresses noise and barren plateaus
\end{itemize}

\medskip

\textbf{Weaknesses}
\begin{itemize}
    \item Needs explicit cost functions, or links to article
    \item Insufficient comparison with existing methods
    \item Could perhaps add complexity analysis or refer to article
\end{itemize}

\medskip


\end{frame}

\begin{frame}{Chapter 5: Subsystem Computation}
\textbf{Strengths}
\begin{itemize}
    \item Highly relevant for NISQ devices
    \item Strong conceptual direction
\end{itemize}

\medskip

\textbf{Weaknesses}
\begin{itemize}
    \item Concept not fully formalized
    \item Perhaps add scaling and error analysis
    \item Add  connection to earlier chapters
\end{itemize}

\medskip



Clarify subsystem approximation strategy.
\end{frame}

%================================================
\section{Main Improvements}
%================================================

\begin{frame}{Key Improvements}

\begin{enumerate}
    \item Strengthen global narrative
    \item Expand key derivations or link to papers
    \item Add scaling and complexity analysis
    \item Connect chapters explicitly
    \item Add critical discussion sections
\end{enumerate}

\medskip

Core transformation:
\[
\text{Exposition} \rightarrow \text{Scientific synthesis}
\]
\end{frame}

%================================================
\section{Thesis Summary (Strong Version)}
%================================================

\begin{frame}{Thesis Summary: Scientific Problem}
This thesis addresses a central challenge in quantum computing:

\medskip

\[
\text{How do we bridge the gap between physical qubits and useful computation in the NISQ era?}
\]

\medskip

This requires solving four interconnected problems:

\begin{itemize}
    \item Physical realization of qubits
    \item Control and entanglement generation
    \item Efficient circuit compilation
    \item Scalable use of small quantum devices
\end{itemize}
\end{frame}

\begin{frame}{Thesis Contributions}
This thesis makes contributions along the full stack:

\medskip

\textbf{(1) Quantum Hardware}
\begin{itemize}
    \item Modeling entanglement via Coulomb interactions
    \item Analysis of electrons-on-helium platform
\end{itemize}

\medskip

\textbf{(2) Quantum Control}
\begin{itemize}
    \item Design of high-fidelity two-qubit gates
    \item Analysis of noise and control constraints
\end{itemize}
\end{frame}

\begin{frame}{Thesis Contributions (cont.)}
\textbf{(3) Quantum Compiling}
\begin{itemize}
    \item Development of recursive variational quantum compiling
    \item Mitigation of vanishing gradients
\end{itemize}

\medskip

\textbf{(4) Subsystem-Based Computation}
\begin{itemize}
    \item Framework for using small quantum devices
    \item Application to machine learning and energy estimation
\end{itemize}
\end{frame}

\begin{frame}{Conceptual Contribution}
The central conceptual contribution is:

\[
\boxed{
\text{A unified framework connecting hardware, control, compilation, and computation}
}
\]

\medskip

This is particularly relevant for:
\begin{itemize}
    \item NISQ devices
    \item noisy quantum algorithms
    \item scalable hybrid approaches
\end{itemize}
\end{frame}

\begin{frame}{Scientific Impact}
The thesis contributes to:

\begin{itemize}
    \item Understanding of physically realizable qubit systems
    \item Development of noise-resilient compiling methods
    \item Strategies for extending small quantum devices
\end{itemize}

\medskip

\[
\text{Impact spans hardware, algorithms, and computational paradigms}
\]
\end{frame}

\begin{frame}{Limitations and Outlook}
Key limitations:

\begin{itemize}
    \item Scalability of proposed methods
    \item Noise sensitivity
    \item Need for experimental validation
\end{itemize}

\medskip

Future directions:

\begin{itemize}
    \item Integration with quantum error mitigation
    \item Hybrid classical-quantum optimization
    \item Larger-scale subsystem architectures
\end{itemize}
\end{frame}

\begin{frame}{Final Assessment}
\[
\boxed{
\text{Strong PhD-level scientific work with high potential impact}
}
\]

\medskip

Main remaining step:

\[
\boxed{
\text{Transform chapters into a coherent scientific narrative}
}
\]


\end{frame}

\end{document}
